% Options for packages loaded elsewhere
\PassOptionsToPackage{unicode}{hyperref}
\PassOptionsToPackage{hyphens}{url}
\PassOptionsToPackage{dvipsnames,svgnames,x11names}{xcolor}
%
\documentclass[
  11pt,
]{article}
\usepackage{amsmath,amssymb}
\usepackage{iftex}
\ifPDFTeX
  \usepackage[T1]{fontenc}
  \usepackage[utf8]{inputenc}
  \usepackage{textcomp} % provide euro and other symbols
\else % if luatex or xetex
  \usepackage{unicode-math} % this also loads fontspec
  \defaultfontfeatures{Scale=MatchLowercase}
  \defaultfontfeatures[\rmfamily]{Ligatures=TeX,Scale=1}
\fi
\usepackage{lmodern}
\ifPDFTeX\else
  % xetex/luatex font selection
\fi
% Use upquote if available, for straight quotes in verbatim environments
\IfFileExists{upquote.sty}{\usepackage{upquote}}{}
\IfFileExists{microtype.sty}{% use microtype if available
  \usepackage[]{microtype}
  \UseMicrotypeSet[protrusion]{basicmath} % disable protrusion for tt fonts
}{}
\makeatletter
\@ifundefined{KOMAClassName}{% if non-KOMA class
  \IfFileExists{parskip.sty}{%
    \usepackage{parskip}
  }{% else
    \setlength{\parindent}{0pt}
    \setlength{\parskip}{6pt plus 2pt minus 1pt}}
}{% if KOMA class
  \KOMAoptions{parskip=half}}
\makeatother
\usepackage{xcolor}
\usepackage[margin=1in]{geometry}
\usepackage{listings}
\newcommand{\passthrough}[1]{#1}
\lstset{defaultdialect=[5.3]Lua}
\lstset{defaultdialect=[x86masm]Assembler}
\setlength{\emergencystretch}{3em} % prevent overfull lines
\providecommand{\tightlist}{%
  \setlength{\itemsep}{0pt}\setlength{\parskip}{0pt}}
\setcounter{secnumdepth}{-\maxdimen} % remove section numbering
\providecommand{\implies}{\Longrightarrow}
\usepackage{listings}
\usepackage{xcolor}
\lstset{breaklines=true, breakatwhitespace=false, basicstyle=\ttfamily\small, columns=fullflexible, postbreak=\mbox{\textcolor{gray}{$\hookrightarrow$}\space}}
\ifLuaTeX
  \usepackage{selnolig}  % disable illegal ligatures
\fi
\IfFileExists{bookmark.sty}{\usepackage{bookmark}}{\usepackage{hyperref}}
\IfFileExists{xurl.sty}{\usepackage{xurl}}{} % add URL line breaks if available
\urlstyle{same}
\hypersetup{
  pdftitle={Exercises from Books: LTI ODE with Impulsive Loads},
  pdfauthor={Denis Pleshkov ~},
  colorlinks=true,
  linkcolor={blue},
  filecolor={Maroon},
  citecolor={blue},
  urlcolor={blue},
  pdfcreator={LaTeX via pandoc}}

\title{Exercises from Books: LTI ODE with Impulsive Loads}
\author{Denis Pleshkov
~\href{mailto:std.approach@gmail.com}{std.approach@gmail.com}}
\date{Last Modified: August 25, 2026}

\begin{document}
\maketitle

{
\hypersetup{linkcolor=}
\setcounter{tocdepth}{3}
\tableofcontents
}
\hypertarget{abstract}{%
\section{Abstract}\label{abstract}}

This article compiles and independently verifies a curated benchmark set
of worked examples, drawn from published textbooks, illustrating the
solution of linear time-invariant (LTI) ordinary differential equations
(ODEs) forced by the Dirac delta function and its derivatives --- the
standard mathematical model of an impulsive load. No new mathematical
results are derived here: each of the thirty-seven solved exercises in
the ``Solved Exercises'' section reproduces a problem statement and its
published closed-form solution essentially verbatim from its source,
with full attribution by author, page, and edition. What distinguishes
this compilation from a plain survey is that every one of those
thirty-seven solutions was additionally checked, by symbolic
Laplace-transform matching, segment-wise ODE and jump-matching, or
numeric recomputation, against its own stated equation and initial
conditions; the handful of transcription issues this uncovered ---
including one textbook example whose printed answer did not actually
satisfy its own stated problem --- are corrected inline with an explicit
derivation rather than left silently as-is. Each solved example
additionally carries a ready-to-run Wolfram Language snippet that a
reader can paste directly into WolframAlpha to reproduce its closed-form
solution independently. A machine-readable export (JSON and CSV) of the
full verified example set accompanies the article. A further set of
unsolved exercise references, spanning twenty additional textbooks, is
indexed to guide further practice. The compilation serves two purposes.
First, it is pedagogical: gathering material otherwise scattered across
dozens of engineering-mathematics, vibrations, and control-theory
textbooks into one alphabetically organized reference, suitable as a
starting point for a reader studying impulsively forced LTI systems.
Second, it is practical: because each solved exercise pairs a fully
specified, independently verified ODE with a known analytical solution,
the collection functions directly as a benchmark/regression-test suite
for validating symbolic or numerical ODE-solving software. The
Verification Methodology and Conclusion sections detail the checking
process and summarize the scope, composition, and limitations of the
collected material.

\hypertarget{keywords}{%
\section{Keywords}\label{keywords}}

Dirac delta function, impulse response, linear time-invariant ODE,
initial value problem, Laplace transform, textbook exercise compilation,
benchmark test suite for ODE solvers, symbolic verification,
WolframAlpha, Wolfram Language

\hypertarget{introduction}{%
\section{Introduction}\label{introduction}}

This article is \textbf{pedagogical} in purpose. It is intended as a
starting point for a reader who wants to study, in a structured way,
linear time-invariant (LTI) ordinary differential equations subject to
impulsive loads (the Dirac delta function and its derivatives as a
forcing term). The article itself contains \textbf{no new mathematical
results}: every worked example in the ``Solved Exercises'' section is
extracted, essentially verbatim, from an existing textbook. Its only
contribution is the \emph{gathering, attribution, and organization} of
material that is otherwise scattered across dozens of books on
differential equations, vibrations, and control theory.

A second, practical motivation for this compilation is \textbf{software
testing}. Because each solved exercise below pairs a well-defined LTI
ODE (with explicit initial conditions) with a known closed-form
analytical solution, the collection as a whole forms a ready-made
benchmark suite. It can be used to validate an existing symbolic or
numerical ODE-solving library, or to build a moderate-sized
regression/unit-test suite for a library still under development ---
simply compare the library's output against the analytical solution
quoted here. Unlike a raw transcription, however, a benchmark is only as
trustworthy as its answer key: every closed-form solution in the Solved
Exercises section was therefore independently checked against its own
stated equation and initial conditions before being included (see the
Verification Methodology section), and the full, verified set is
additionally provided as a machine-readable JSON/CSV export for direct
use in an automated test harness. For a quicker, example-by-example spot
check, every solved exercise also carries its own Wolfram Language
snippet that can be pasted directly into WolframAlpha to reproduce that
one closed-form solution on the spot, without setting up a full test
harness.

The material is split into two categories:

\begin{enumerate}
\def\labelenumi{\arabic{enumi}.}
\tightlist
\item
  \textbf{Solved Exercises} --- problems taken directly from a textbook
  together with their published closed-form solution. Nothing here is
  derived by the present author; it is a curated extraction of existing
  results, organized by source and alphabetized by author.
\item
  \textbf{Additional Exercises} --- pointers to further exercises in the
  same (and other) textbooks that are relevant to impulsive loading but
  for which no solution is reproduced here. These are
  page/problem-number references only, meant to guide further practice
  and reading.
\end{enumerate}

Both categories are ordered alphabetically by author surname, and the
bibliography follows the same convention.

\begin{center}\rule{0.5\linewidth}{0.5pt}\end{center}

\hypertarget{solved-exercises}{%
\section{Solved Exercises}\label{solved-exercises}}

Worked examples reproduced from the source textbooks, including the
original problem statement and its published closed-form solution.
Sorted alphabetically by (first) author. Following academic convention,
verbatim text taken from a source (an example's title or its problem
statement) is set in double quotation marks; page/example locators and
the present author's own classification notes are left unquoted.

Each entry also carries a bracketed \emph{WolframAlpha verification}
block: a Wolfram Language command that can be pasted into the input box
at \href{https://www.wolframalpha.com}{wolframalpha.com}, which returns
the result independently. Most of these commands were constructed by
hand and cross-checked symbolically offline rather than run live (see
Verification Methodology below), so pasting one in may be giving it its
first live test.

\hypertarget{boyce-diprima-elementary-differential-equations-and-boundary-value-problems}{%
\subsection{Boyce \& DiPrima --- Elementary differential equations and
boundary value
problems}\label{boyce-diprima-elementary-differential-equations-and-boundary-value-problems}}

p.272 Initial Value Problem (IVP) for 2nd order with zero Initial
Condition (IC) \[
\begin{aligned}
&2y'' + y' + 2y = \delta(t - 5) \\
&y(0) = 0, \quad y'(0) = 0
\end{aligned}
\implies y(t) =
\begin{cases}
0, & t < 5 \\[6pt]
\dfrac{2}{\sqrt{15}} e^{-(t-5)/4} \sin\left(\dfrac{\sqrt{15}}{4}(t-5)\right), & t \ge 5
\end{cases}
\] \emph{{[}WolframAlpha:}

\begin{lstlisting}
2 y''[t] + y'[t] + 2 y[t] == Delta[t - 5], y[0] == 0, y'[0] == 0
\end{lstlisting}

\emph{Returns
\(y(t) = \frac{2 \, e^{5/4 - t/4} \, u(t-5) \, \sin\left(\frac{1}{4} \, \sqrt{15} \, (t-5)\right)}{\sqrt{15}}\)}{]}

\emph{{[}Jump condition at \(t=5\): the equation
\(2y''+y'+2y=\delta(t-5)\) normalizes to
\(y''+\tfrac12 y'+y=\tfrac12\delta(t-5)\). Phase vector
\(\mathbf y=(y,y')\) changes by
\(\Delta\mathbf y(5)=\left(0,\tfrac12\right)\).{]}}

\hypertarget{campbell-haberman-introduction-to-differential-equations-with-dynamical-systems}{%
\subsection{Campbell \& Haberman --- Introduction to differential
equations with dynamical
systems}\label{campbell-haberman-introduction-to-differential-equations-with-dynamical-systems}}

p.263 IVP for 1st order system
\[ y' + y = \delta(t - 1), \, y(0) = 1 \implies  y(t)=\begin{cases}
e^{-t}, & t < 1 \\[4pt]
e^{-t} + e^{-(t-1)}, & t \ge 1
\end{cases}\] \emph{{[}WolframAlpha:}

\begin{lstlisting}
y'[t] + y[t] == Delta[t - 1], y[0] == 1
\end{lstlisting}

\emph{Returns \(y(t) = e^{-t} + e^{-(t-1)} \, u(t-1)\)}{]}

\emph{{[}Jump condition at \(t=1\): \(y'+y=\delta(t-1)\) is first order.
Phase vector \(\mathbf y=(y)\) changes by
\(\Delta\mathbf y(1)=(1)\).{]}}

\hypertarget{edwards-penney-elementary-differential-equations-with-boundary-value-problems}{%
\subsection{Edwards \& Penney --- Elementary differential equations with
boundary value
problems}\label{edwards-penney-elementary-differential-equations-with-boundary-value-problems}}

p.318 The IVP is \[
  x'' + 4x = 8\delta_{2\pi}(t); \; x(0) = 3, \; x'(0) = 0 \implies x(t)= \begin{cases}
3\cos 2t, & t < 2\pi \\[4pt]
3\cos 2t + 4\sin 2t, & t \ge 2\pi
\end{cases}
\] \emph{{[}WolframAlpha:}

\begin{lstlisting}
x''[t] + 4 x[t] == 8 Delta[t - 2 Pi], x[0] == 3, x'[0] == 0
\end{lstlisting}

\emph{Returns \(x(t) = 3\cos(2t) + 4\sin(2t) \, u(t-2\pi)\)}{]}

\emph{{[}Jump condition at \(t=2\pi\): phase vector \(\mathbf x=(x,x')\)
changes by \(\Delta\mathbf x(2\pi)=(0,8)\), the impulse strength.{]}}

\hypertarget{esfandiari-lu-modeling-and-analysis-of-dynamic-systems}{%
\subsection{Esfandiari \& Lu --- Modeling and analysis of dynamic
systems}\label{esfandiari-lu-modeling-and-analysis-of-dynamic-systems}}

p.57 ``Example 2.27: Initial Condition \(\neq\) Initial Value'' \[
 \ddot{x} + \dot{x} + 2x = \delta(t), \, x(0^-) = 0, \, \dot{x}(0^-) = 0
\implies \dot{x}(0^+) =  1
\] \emph{{[}WolframAlpha:}

\begin{lstlisting}
x''[t] + x'[t] + 2 x[t] == Delta[t], x[0] == 0, x'[0] == 0
\end{lstlisting}

\emph{Returns
\(x(t) = \frac{2 \, e^{-t/2} \, u(t) \, \sin\left(\frac{\sqrt{7}}{2} t\right)}{\sqrt{7}}\)
(the book only states the jump \(\dot x(0^+)=1\); this full closed form
is consistent with it --- differentiating gives \(\dot x(0^+)=1\))}{]}

\emph{{[}Jump condition at \(t=0\): confirming the book's own claim,
\(\ddot x+\dot x+2x=\delta(t)\). Phase vector \(\mathbf x=(x,\dot x)\)
changes by \(\Delta\mathbf x(0)=(0,1)\).{]}}

p.343 ``Impulse Response of First-Order Systems''
\[ x(t) = e^{-t/\tau} x_0 + \frac{A}{\tau} e^{-t/\tau} \] where A -
impulse's magnitude \(\tau\) - coefficient for higher derivative in ODE.
\emph{{[}WolframAlpha:}

\begin{lstlisting}
tau x'[t] + x[t] == A Delta[t], x[0] == x0
\end{lstlisting}

\emph{Returns
\(x(t) = e^{-t/\tau} x_0 + \frac{A}{\tau} e^{-t/\tau}\)}{]}

\emph{{[}Jump condition at \(t=0\): normalizing \(\tau x'+x=A\delta(t)\)
to \(x'+\tfrac1\tau x=\tfrac{A}{\tau}\delta(t)\) (first order). Phase
vector \(\mathbf x=(x)\) changes by
\(\Delta\mathbf x(0)=\left(\dfrac{A}{\tau}\right)\).{]}}

\hypertarget{franklin-powell-emami-naeini-feedback-control-of-dynamic-systems}{%
\subsection{Franklin, Powell \& Emami-Naeini --- Feedback control of
dynamic
systems}\label{franklin-powell-emami-naeini-feedback-control-of-dynamic-systems}}

p.~110 \[
\dot{y} + ky = u = \delta(t), y(0) = 0 \equiv \dot{y} + ky = 0, \quad y(0^+) = 1
\] \emph{{[}WolframAlpha:}

\begin{lstlisting}
y'[t] + k y[t] == Delta[t], y[0] == 0
\end{lstlisting}

\emph{Returns \(y(t) = e^{-k t} \, u(t)\) (at \(t=0^+\) this gives
\(y(0^+)=1\), matching the book)}{]}

\emph{{[}Jump condition at \(t=0\): confirming the book's own reduction,
\(y'+ky=\delta(t)\) is first order. Phase vector \(\mathbf y=(y)\)
changes by \(\Delta\mathbf y(0)=(1)\).{]}}

p.151 \[
H(s) = \frac{2s + 1}{(s+1)^2 + 2^2} \implies h(t) = \left( 2e^{-t} \cos 2t - \frac{1}{2}e^{-t} \sin 2t \right) 1(t)
\] \emph{{[}WolframAlpha:}

\begin{lstlisting}
InverseLaplaceTransform[(2 s + 1)/((s + 1)^2 + 4), s, t]
\end{lstlisting}

\emph{Returns
\(h(t) = 2 e^{-t} \cos(2t) - \frac{1}{2} e^{-t} \sin(2t)\)}{]}

\emph{{[}Jump condition at \(t=0\): because the right-hand side carries
a derivative of the (impulsive) input, \(2x'+x\) with \(x=\delta(t)\),
the relative degree of \(H(s)\) drops to \(1\). Phase vector
\(\mathbf y=(y,y')\) changes by \(\Delta\mathbf y(0)=(2,-3)\).{]}}

\hypertarget{gangadharaiah-sandeep-engineering-applications-of-the-laplace-transform}{%
\subsection{Gangadharaiah \& Sandeep --- Engineering applications of the
Laplace
transform}\label{gangadharaiah-sandeep-engineering-applications-of-the-laplace-transform}}

p.239 Example 3.8. ``Use the Laplace transform to find \ldots{} the
impulse response of the system if the differential equation describes
the system'' \[
\frac{d^2y(t)}{dt^2} + 5\frac{dy(t)}{dt} + 6y(t) = \frac{d^2x(t)}{dt^2} + 8\frac{dx(t)}{dt} + 13x(t) \implies h(t) = \delta(t) + e^{-2t} + 2e^{-3t}
\] \emph{{[}WolframAlpha:}

\begin{lstlisting}
InverseLaplaceTransform[(s^2 + 8 s + 13)/(s^2 + 5 s + 6), s, t]
\end{lstlisting}

\emph{Returns \(h(t) = \delta(t) + e^{-2t} + 2 e^{-3t}\)}{]}

\emph{{[}Jump condition at \(t=0\): numerator and denominator of
\(H(s)=(s^2+8s+13)/(s^2+5s+6)\) share degree \(2\), so the input's own
\(\delta(t)\) passes straight through as a direct-feedthrough term
(coefficient \(=1\)) on top of the smooth part
\(h_{\text{reg}}(t)=e^{-2t}+2e^{-3t}\). Phase vector
\(\mathbf h_{\text{reg}}=(h_{\text{reg}},h_{\text{reg}}')\) changes by
\(\Delta\mathbf h_{\text{reg}}(0)=(3,-8)\).{]}}

p.248 Example 3.11. ``Consider the causal LTI system described by the
second differential equation \ldots{} Determine the impulse response''
\[
\frac{d^2y(t)}{dt^2} + 5\frac{dy(t)}{dt} + 6y(t) = x(t) \implies h(t) = e^{-2t} - e^{-3t}.
\] \emph{{[}WolframAlpha:}

\begin{lstlisting}
y''[t] + 5 y'[t] + 6 y[t] == Delta[t], y[0] == 0, y'[0] == 0
\end{lstlisting}

\emph{Returns \(h(t) = e^{-2t} - e^{-3t}\)}{]}

\emph{{[}Jump condition at \(t=0\): \(y''+5y'+6y=\delta(t)\). Phase
vector \(\mathbf y=(y,y')\) changes by \(\Delta\mathbf y(0)=(0,1)\).{]}}

p.254 Example 3.12. ``find the impulse response of the system if the
third-order differential equation describes the system'' \[
\frac{d^3 y(t)}{dt^3} + 6 \frac{d^2 y(t)}{dt^2} + 11 \frac{dy(t)}{dt} + 6y(t) = x(t) \implies h(t) = \frac{1}{2} e^{-t} - e^{-2t} + \frac{1}{2} e^{-3t}.
\] \emph{{[}WolframAlpha:}

\begin{lstlisting}
y'''[t] + 6 y''[t] + 11 y'[t] + 6 y[t] == Delta[t], y[0] == 0, y'[0] == 0, y''[0] == 0
\end{lstlisting}

\emph{Returns
\(h(t) = \frac{1}{2} e^{-t} - e^{-2t} + \frac{1}{2} e^{-3t}\)}{]}

\emph{{[}Jump condition at \(t=0\): this third-order equation. Phase
vector \(\mathbf y=(y,y',y'')\) changes by
\(\Delta\mathbf y(0)=(0,0,1)\).{]}}

p.257 Example 3.14. ``Compute the impulse response of the transform with
the transfer function'' \[
H(s) = \frac{s^2 - s + 1}{s^2 + 2s + 1} \implies y(t) = \delta(t) - 3e^{-t} + 3te^{-t}.
\] \emph{{[}WolframAlpha:}

\begin{lstlisting}
InverseLaplaceTransform[(s^2 - s + 1)/(s^2 + 2 s + 1), s, t]
\end{lstlisting}

\emph{Returns \(y(t) = \delta(t) - 3 e^{-t} + 3 t e^{-t}\)}{]}

\emph{{[}Jump condition at \(t=0\): as in Example 3.8,
\(\deg N=\deg D=2\) for \(H(s)=(s^2-s+1)/(s^2+2s+1)\), giving a direct
feedthrough of coefficient \(1\) plus a smooth remainder
\(y_{\text{reg}}(t)=-3e^{-t}+3te^{-t}\). Phase vector
\(\mathbf y_{\text{reg}}=(y_{\text{reg}},y_{\text{reg}}')\) changes by
\(\Delta\mathbf y_{\text{reg}}(0)=(-3,6)\).{]}}

p.288 Example 3.27. ``find impulse response of the system if the
differential equation describes the system'' \[
\frac{d^2z(t)}{dt^2} + 3\frac{dz(t)}{dt} + 2z(t) = \frac{d^2x(t)}{dt^2} + 6\frac{dx(t)}{dt} + 7x(t) \implies h(t) = \delta(t) + 2e^{-t} + e^{-2t}.
\] \emph{{[}WolframAlpha:}

\begin{lstlisting}
InverseLaplaceTransform[(s^2 + 6 s + 7)/(s^2 + 3 s + 2), s, t]
\end{lstlisting}

\emph{Returns \(h(t) = \delta(t) + 2 e^{-t} + e^{-2t}\)}{]}

\emph{{[}Jump condition at \(t=0\): again \(\deg N=\deg D=2\) for
\(H(s)=(s^2+6s+7)/(s^2+3s+2)\), giving a direct feedthrough of
coefficient \(1\) plus a smooth remainder
\(h_{\text{reg}}(t)=2e^{-t}+e^{-2t}\). Phase vector
\(\mathbf h_{\text{reg}}=(h_{\text{reg}},h_{\text{reg}}')\) changes by
\(\Delta\mathbf h_{\text{reg}}(0)=(3,-4)\).{]}}

p.289 Example 3.28. ``find the impulse response of the system if the
differential equation describes the system'' \[
\frac{d^2z(t)}{dt^2} + 4\frac{dz(t)}{dt} + 10z(t) = x(t) \implies h(t) = \frac{1}{\sqrt{6}} e^{-2t} \sin\left(\sqrt{6}t\right).
\] \emph{{[}WolframAlpha:}

\begin{lstlisting}
z''[t] + 4 z'[t] + 10 z[t] == Delta[t], z[0] == 0, z'[0] == 0
\end{lstlisting}

\emph{Returns
\(h(t) = \frac{1}{\sqrt{6}} e^{-2t} \sin\left(\sqrt{6} \, t\right)\)}{]}

\emph{{[}Jump condition at \(t=0\): \(z''+4z'+10z=\delta(t)\). Phase
vector \(\mathbf z=(z,z')\) changes by \(\Delta\mathbf z(0)=(0,1)\).{]}}

p.306 Example 4.7. ``Solve the initial value problem'' \[
\frac{d^2 y(t)}{dt^2} + 5 \frac{dy(t)}{dt} + 6y(t) = \delta(t - \pi) - \delta(t - 2\pi)
\]

with \(y(0) = 0 = y'(0)\). \[
y(t) = \left( e^{-2(t-\pi)} - e^{-3(t-\pi)} \right) u(t - \pi) - \left( e^{-2(t-2\pi)} - e^{-3(t-2\pi)} \right) u(t - 2\pi).
\] \emph{{[}WolframAlpha:}

\begin{lstlisting}
y''[t] + 5 y'[t] + 6 y[t] == Delta[t - Pi] - Delta[t - 2 Pi], y[0] == 0, y'[0] == 0
\end{lstlisting}

\emph{Returns
\(y(t) = \left(e^{-2(t-\pi)} - e^{-3(t-\pi)}\right) u(t-\pi) - \left(e^{-2(t-2\pi)} - e^{-3(t-2\pi)}\right) u(t-2\pi)\)}{]}

\emph{{[}Jump conditions: phase vector \(\mathbf y=(y,y')\) changes by
\(\Delta\mathbf y(\pi)=(0,1)\) at \(t=\pi\), and by
\(\Delta\mathbf y(2\pi)=(0,-1)\) at \(t=2\pi\) (the coefficient of
\(-\delta(t-2\pi)\)).{]}}

p.340 Example 4.25. ``Obtain the solution of the fourth-order
differential equation'' \[
\frac{d^4 y(t)}{dt^4} + 2 \frac{d^3 y(t)}{dt^3} - \frac{d^2 y(t)}{dt^2} - 2 \frac{dy(t)}{dt} = \delta(t)
\]

along with the initial condition \[
y(0) = 1 \quad \text{and} \quad y'(0) = y''(0) = y'''(0) = 0.
\]

\[
y(t) = \frac{1}{2} + \frac{1}{2}e^{-t} + \frac{1}{6}e^{t} - \frac{1}{6}e^{-2t}.
\] \emph{{[}WolframAlpha:}

\begin{lstlisting}
y''''[t] + 2 y'''[t] - y''[t] - 2 y'[t] == Delta[t], y[0] == 1, y'[0] == 0, y''[0] == 0, y'''[0] == 0
\end{lstlisting}

\emph{Returns
\(y(t) = \frac{1}{2} + \frac{1}{2} e^{-t} + \frac{1}{6} e^{t} - \frac{1}{6} e^{-2t}\)}{]}

\emph{{[}Jump condition at \(t=0\): this fourth-order equation. Phase
vector \(\mathbf y=(y,y',y'',y''')\) changes by
\(\Delta\mathbf y(0)=(0,0,0,1)\).{]}} \emph{{[}Editorial note: as
transcribed, this example's stated answer
(\(y = \frac{1}{2} - e^t + \frac{3}{2}e^{2t}\)) does not satisfy its own
stated differential equation and initial conditions --- it fails the
homogeneous-equation check for \(t>0\) and the required continuity of
\(y'\) and \(y''\) at \(t=0\). The closed-form solution above is the
unique function consistent with the stated fourth-order equation,
\(y(0)=1\), \(y'(0)=y''(0)=y'''(0)=0\), and \(\delta(t)\) forcing; it
was re-derived via the Laplace transform and independently confirmed by
direct substitution back into the differential equation. It replaces the
original transcription here as a high-confidence, mathematically
necessary correction rather than a silent guess.{]}}

p.387 Example 4.49. ``Obtain the solution of the second-order
differential equation'' \[
\frac{d^2y(t)}{dt^2} + 5\frac{dy(t)}{dt} + 6y(t) = 3\delta(t-2) - 4\delta(t-4)
\]

along with the initial conditions \(y(0) = 0 = y'(0)\). \[
y(t) = 3(e^{-2(t-2)} - e^{-3(t-2)})u(t-2) - 4(e^{-2(t-4)} - e^{-3(t-4)})u(t-4).
\] \emph{{[}WolframAlpha:}

\begin{lstlisting}
y''[t] + 5 y'[t] + 6 y[t] == 3 Delta[t - 2] - 4 Delta[t - 4], y[0] == 0, y'[0] == 0
\end{lstlisting}

\emph{Returns
\(y(t) = 3\left(e^{-2(t-2)} - e^{-3(t-2)}\right) u(t-2) - 4\left(e^{-2(t-4)} - e^{-3(t-4)}\right) u(t-4)\)}{]}

\emph{{[}Jump conditions: phase vector \(\mathbf y=(y,y')\) changes by
\(\Delta\mathbf y(2)=(0,3)\) at \(t=2\) (the coefficient of
\(3\delta(t-2)\)), and by \(\Delta\mathbf y(4)=(0,-4)\) at \(t=4\) (the
coefficient of \(-4\delta(t-4)\)).{]}}

\hypertarget{lathi-green-linear-systems-and-signals}{%
\subsection{Lathi \& Green --- Linear systems and
signals}\label{lathi-green-linear-systems-and-signals}}

p.164 ``EXAMPLE 2.5 Impulse Response via Impulse Matching'' ``Find the
impulse response h(t) for a system specified by (D2 +5D+6)y(t) =
(D+1)x(t)'' Solution \[ h(t) = (-e^{-2t} + 2e^{-3t})u(t) \]
\emph{{[}WolframAlpha:}

\begin{lstlisting}
InverseLaplaceTransform[(s + 1)/(s^2 + 5 s + 6), s, t]
\end{lstlisting}

\emph{Returns \(h(t) = \left(-e^{-2t} + 2 e^{-3t}\right) u(t)\)}{]}

\emph{{[}Jump condition at \(t=0\): because the input enters as
\((D+1)x\), the relative degree of the transfer function is \(1\). Phase
vector \(\mathbf y=(y,y')\) changes by
\(\Delta\mathbf y(0)=(1,-4)\).{]}}

p.166 EXAMPLE 2.6 ``Determine the unit impulse response h(t) for a
system specified by the equation (D\^{}2 +3D+2)y(t) = Dx(t)'' Solution
\[ h(t) = (-e^{-t} + 2e^{-2t})u(t) \] \emph{{[}WolframAlpha:}

\begin{lstlisting}
InverseLaplaceTransform[s/(s^2 + 3 s + 2), s, t]
\end{lstlisting}

\emph{Returns \(h(t) = \left(-e^{-t} + 2 e^{-2t}\right) u(t)\)}{]}

\emph{{[}Jump condition at \(t=0\): the input enters as \(Dx\), again
giving relative degree \(1\). Phase vector \(\mathbf y=(y,y')\) changes
by \(\Delta\mathbf y(0)=(1,-3)\).{]}}

p.167 ``DRILL 2.4 Finding the Impulse Response'' ``Determine the unit
impulse response of LTIC systems described by the following equations:''

\begin{enumerate}
\def\labelenumi{(\alph{enumi})}
\item
  \((D + 2)y(t) = (3D + 5)x(t) \implies h(t)=3\delta(t) - e^{-2t}u(t)\)
\item
  \(D(D + 2)y(t) = (D + 4)x(t) \implies h(t)=(2 - e^{-2t})u(t)\)
\item
  \((D^2 + 2D + 1)y(t) = Dx(t) \implies h(t)=(1 - t)e^{-t}u(t)\)
  \emph{{[}WolframAlpha:}
\end{enumerate}

\begin{lstlisting}
InverseLaplaceTransform[(3 s + 5)/(s + 2), s, t]
InverseLaplaceTransform[(s + 4)/(s (s + 2)), s, t]
InverseLaplaceTransform[s/(s + 1)^2, s, t]
\end{lstlisting}

\emph{Returns, respectively: \(h(t) = 3\delta(t) - e^{-2t} u(t)\),
\(\; h(t) = \left(2 - e^{-2t}\right) u(t)\),
\(\; h(t) = (1-t) e^{-t} u(t)\)}{]}

\emph{{[}Jump conditions at \(t=0\), one per part: (a)
\((D+2)y=(3D+5)x\) has \(\deg N=\deg D=1\), a direct feedthrough of
coefficient \(3\) plus a remainder that changes by
\(\Delta\mathbf y_{\text{reg}}(0)=(-1)\); (b) \(D(D+2)y=(D+4)x\) has
relative degree \(1\): \(\Delta\mathbf y(0)=(1,2)\); (c)
\((D^2+2D+1)y=Dx\) likewise: \(\Delta\mathbf y(0)=(1,-2)\).{]}}

\hypertarget{nagle-saff-snider-fundamentals-of-differential-equations}{%
\subsection{Nagle, Saff \& Snider --- Fundamentals of differential
equations}\label{nagle-saff-snider-fundamentals-of-differential-equations}}

p.403 Example 4 ``A linear system is governed by the differential
equation'' \[
y'' + 2y' + 5y = \delta(t), \quad y(0) = 0, \quad y'(0) = 0 \implies y(t) = \frac{1}{2} e^{-t} \sin 2t
\] \emph{{[}WolframAlpha:}

\begin{lstlisting}
y''[t] + 2 y'[t] + 5 y[t] == Delta[t], y[0] == 0, y'[0] == 0
\end{lstlisting}

\emph{Returns \(y(t) = \frac{1}{2} e^{-t} \sin(2t)\)}{]}

\emph{{[}Jump condition at \(t=0\): phase vector \(\mathbf y=(y,y')\)
changes by \(\Delta\mathbf y(0)=(0,1)\).{]}} \emph{{[}Editorial note: as
transcribed, this equation was incomplete --- missing the forcing term
on the left-hand side and the output variable on the right (the
``\(\implies =\)'' read as a fragment). The equation above,
\(y''+2y'+5y=\delta(t)\) with zero initial conditions, is the unique
standard-form IVP whose closed-form solution matches the book's own
stated answer exactly (characteristic roots \(-1\pm2i\) give precisely
\(\frac{1}{2}e^{-t}\sin2t\)), and is used here as a high-confidence
reconstruction rather than left incomplete.{]}}

p.409 Example 1 \[
\frac{d^2x}{dt^2} + 9x = 3\delta(t - \pi); \quad x(0) = 1, \quad \frac{dx}{dt}(0) = 0 \implies x(t) =
\begin{cases}
\cos 3t, & t < \pi, \\
\cos 3t - \sin 3t, & \pi < t
\end{cases}
\] \emph{{[}WolframAlpha:}

\begin{lstlisting}
x''[t] + 9 x[t] == 3 Delta[t - Pi], x[0] == 1, x'[0] == 0
\end{lstlisting}

\emph{Returns \(x(t) = \cos(3t) - \sin(3t) \, u(t-\pi)\)}{]}

\emph{{[}Jump condition at \(t=\pi\): phase vector \(\mathbf x=(x,x')\)
changes by \(\Delta\mathbf x(\pi)=(0,3)\), the impulse strength.{]}}

\hypertarget{nagy-ordinary-differential-equations}{%
\subsection{Nagy --- Ordinary differential
equations}\label{nagy-ordinary-differential-equations}}

p.202 \[
y'' + \omega_0^2 y = f_0 \delta(t - t_0), \quad y(0) = y_0, \quad y'(0) = 0 \implies y(t) = y_0 \cos(\omega_0 t) + \frac{f_0}{\omega_0} u(t - t_0) \sin(\omega_0 (t - t_0))
\] \emph{{[}WolframAlpha:}

\begin{lstlisting}
y''[t] + w0^2 y[t] == f0 Delta[t - t0], y[0] == y0, y'[0] == 0
\end{lstlisting}

\emph{Returns
\(y(t) = y_0 \cos(\omega_0 t) + \frac{f_0}{\omega_0} \, u(t-t_0) \sin\left(\omega_0 (t-t_0)\right)\)}{]}

\emph{{[}Jump condition at \(t=t_0\): phase vector \(\mathbf y=(y,y')\)
changes by \(\Delta\mathbf y(t_0)=(0,f_0)\), the impulse strength.{]}}

p.205 Example 3.4.6. ``Find the impulse response function'' \[
L(y) = y'' + 2y' + 2y \implies y_\delta(t) = u(t-c)e^{-(t-c)}\sin(t-c)
\] \emph{{[}WolframAlpha:}

\begin{lstlisting}
y''[t] + 2 y'[t] + 2 y[t] == Delta[t], y[0] == 0, y'[0] == 0
\end{lstlisting}

\emph{Returns \(y_\delta(t) = e^{-t} \sin(t) \, u(t)\) (the \(c=0\) case
of the book's general \(u(t-c)e^{-(t-c)}\sin(t-c)\))}{]}

\emph{{[}Jump condition at \(t=c\): phase vector \(\mathbf y=(y,y')\)
changes by \(\Delta\mathbf y(c)=(0,1)\).{]}}

p.205 Example 3.4.7. ``Find the solution y to the initial value
problem'' \[
y'' - y = -20 \delta(t-3), \quad y(0) = 1, \quad y'(0) = 0 \implies y(t) = \cosh(t) - 20 \, u(t-3) \, \sinh(t-3)
\] \emph{{[}WolframAlpha:}

\begin{lstlisting}
y''[t] - y[t] == -20 Delta[t - 3], y[0] == 1, y'[0] == 0
\end{lstlisting}

\emph{Returns \(y(t) = \cosh(t) - 20 \, u(t-3) \sinh(t-3)\)}{]}

\emph{{[}Jump condition at \(t=3\): phase vector \(\mathbf y=(y,y')\)
changes by \(\Delta\mathbf y(3)=(0,-20)\), the impulse strength.{]}}

p.206 Example 3.4.8. ``Find the solution to the initial value problem''
\[
y'' + 4y = \delta(t - \pi) - \delta(t - 2\pi), \quad y(0) = 0, \quad y'(0) = 0 \implies y(t) = \frac{1}{2} \left[ u(t - \pi) - u(t - 2\pi) \right] \sin(2t)
\] \emph{{[}WolframAlpha:}

\begin{lstlisting}
y''[t] + 4 y[t] == Delta[t - Pi] - Delta[t - 2 Pi], y[0] == 0, y'[0] == 0
\end{lstlisting}

\emph{Returns
\(y(t) = \frac{1}{2}\left[u(t-\pi) - u(t-2\pi)\right] \sin(2t)\)}{]}

\emph{{[}Jump conditions: phase vector \(\mathbf y=(y,y')\) changes by
\(\Delta\mathbf y(\pi)=(0,1)\) at \(t=\pi\), and by
\(\Delta\mathbf y(2\pi)=(0,-1)\) at \(t=2\pi\) (the coefficient of
\(-\delta(t-2\pi)\)).{]}}

\hypertarget{ogata-modern-control-engineering}{%
\subsection{Ogata --- Modern control
engineering}\label{ogata-modern-control-engineering}}

p.163 ``Unit-Impulse Response of First-Order Systems'' \[
C(s) = \frac{1}{Ts + 1} \implies c(t) = \frac{1}{T} e^{-t/T}, \quad \text{for } t \geq 0
\] \emph{{[}WolframAlpha:}

\begin{lstlisting}
InverseLaplaceTransform[1/(T s + 1), s, t]
\end{lstlisting}

\emph{Returns \(c(t) = \frac{1}{T} e^{-t/T}\)}{]}

\emph{{[}Jump condition at \(t=0\): normalizing \(Tc'+c=\delta(t)\) to
\(c'+\tfrac1T c=\tfrac1T\delta(t)\) (first order). Phase vector
\(\mathbf c=(c)\) changes by
\(\Delta\mathbf c(0)=\left(\dfrac1T\right)\).{]}}

p.178 ``Impulse Response of Second-Order Systems'' \[
C(s) = \frac{\omega_n^2}{s^2 + 2\zeta\omega_n s + \omega_n^2}
\]

For \(0 \leq \zeta < 1\),

\[
c(t) = \frac{\omega_n}{\sqrt{1 - \zeta^2}} e^{-\zeta \omega_n t} \sin \omega_n \sqrt{1 - \zeta^2} t, \quad \text{for } t \geq 0
\]

For \(\zeta = 1\),

\[
c(t) = \omega_n^2 t e^{-\omega_n t}, \quad \text{for } t \geq 0
\]

For \(\zeta > 1\),

\[
c(t) = \frac{\omega_n}{2 \sqrt{\zeta^2 - 1}} e^{-\left( \zeta - \sqrt{\zeta^2 - 1} \right) \omega_n t} - \frac{\omega_n}{2 \sqrt{\zeta^2 - 1}} e^{-\left( \zeta + \sqrt{\zeta^2 - 1} \right) \omega_n t}, \quad \text{for } t \geq 0
\] \emph{{[}WolframAlpha: this one needs the assumption on \(\zeta\)
made explicit (unlike the other entries, a bare equation list would
leave \(\zeta\)'s range ambiguous), so paste each line separately:}

\begin{lstlisting}
Assuming[0 < zeta < 1 && wn > 0, DSolve[{y''[t] + 2 zeta wn y'[t] + wn^2 y[t] == wn^2 Delta[t], y[0] == 0, y'[0] == 0}, y[t], t]]
Assuming[zeta == 1 && wn > 0, DSolve[{y''[t] + 2 zeta wn y'[t] + wn^2 y[t] == wn^2 Delta[t], y[0] == 0, y'[0] == 0}, y[t], t]]
Assuming[zeta > 1 && wn > 0, DSolve[{y''[t] + 2 zeta wn y'[t] + wn^2 y[t] == wn^2 Delta[t], y[0] == 0, y'[0] == 0}, y[t], t]]
\end{lstlisting}

\emph{Returns, respectively:
\(c(t) = \frac{\omega_n}{\sqrt{1-\zeta^2}} e^{-\zeta\omega_n t}\sin\left(\omega_n\sqrt{1-\zeta^2}\,t\right)\),
\(\; c(t) = \omega_n^2 t \, e^{-\omega_n t}\),
\(\; c(t) = \frac{\omega_n}{2\sqrt{\zeta^2-1}} e^{-(\zeta-\sqrt{\zeta^2-1})\omega_n t} - \frac{\omega_n}{2\sqrt{\zeta^2-1}} e^{-(\zeta+\sqrt{\zeta^2-1})\omega_n t}\)}{]}

\emph{{[}Jump condition at \(t=0\) (all three damping cases): phase
vector \(\mathbf c=(c,c')\) changes by
\(\Delta\mathbf c(0)=(0,\omega_n^2)\), independent of \(\zeta\).{]}}

\hypertarget{shabana-vibration-of-discrete-and-continuous-systems}{%
\subsection{Shabana --- Vibration of discrete and continuous
systems}\label{shabana-vibration-of-discrete-and-continuous-systems}}

p.41 Example 1.10 ``Find the response of the single degree of freedom
system shown in Fig. 17 to the rectangular impulsive force shown in Fig.
16, where m = 10 kg, k = 9,000 N/m, c = 18 N·s/m, and Fo = 10,000 N. The
force is assumed to act at time t = 0 and the impact interval is assumed
to be 0.005 s.'' The system response to the impulsive force is then
given by \[
\begin{aligned}
x(t) &= \frac{l}{m\omega_d} e^{-\xi\omega t} \sin \omega_d t = \frac{50}{(10)(29.986)} e^{-(0.03)(30)t} \sin 29.986t = 0.1667e^{-0.9t} \sin 29.986t
\end{aligned}
\] \emph{{[}WolframAlpha:}

\begin{lstlisting}
m = 10; k = 9000; c = 18; F0 = 10000; dt = 0.005; wn = Sqrt[k/m]; xi = c/(2 Sqrt[m k]); wd = wn Sqrt[1 - xi^2]; N[{wn, xi, wd, xi wn, (F0 dt)/(m wd)}]
\end{lstlisting}

\emph{Returns
\(\{\omega_n,\,\xi,\,\omega_d,\,\xi\omega_n,\,\text{amplitude}\} = \{30,\ 0.03,\ 29.9865,\ 0.9,\ 0.166742\}\)
(matches the book's stated 29.986, 0.9, 0.1667)}{]}

\emph{{[}Jump condition at \(t=0\): the finite-duration force is
idealized as an impulse of magnitude
\(I=F_0\Delta t=(10{,}000)(0.005)=50\ \text{N·s}\) applied to
\(m\ddot x+c\dot x+kx=I\delta(t)\). Phase vector
\(\mathbf x=(x,\dot x)\) changes by
\(\Delta\mathbf x(0)=(0,I/m)=(0,5\ \text{m/s})\), the impulse--momentum
theorem (consistent with the stated amplitude
\(x_0/(m\omega_d)=50/(10\cdot29.986)=0.1667\)).{]}} \emph{{[}Editorial
note: as transcribed, the damping coefficient's unit was garbled (``c =
18 N· slm''). Recomputing \(\omega_n, \xi, \omega_d, \xi\omega_n\), and
the response amplitude from \(m=10\), \(k=9000\), \(c=18\),
\(F_0=10{,}000\), \(\Delta t=0.005\) reproduces every downstream number
the book states (29.986, 0.9, 0.1667) exactly, confirming the numeric
value 18 is correct; only the unit label was corrected here, to the
standard ``N·s/m''.{]}}

\hypertarget{xie-differential-equations-for-engineers}{%
\subsection{Xie --- Differential equations for
engineers}\label{xie-differential-equations-for-engineers}}

p.258 Example 6.11 \[
\mathcal{L}^{-1} \left\{ \frac{s}{(s-2)^5} \right\} \implies f(t) = \frac{1}{12} e^{2t} t^3 (2 + t)
\] \emph{{[}WolframAlpha:}

\begin{lstlisting}
InverseLaplaceTransform[s/(s - 2)^5, s, t]
\end{lstlisting}

\emph{Returns \(f(t) = \frac{1}{12} e^{2t} t^3 (2+t)\)}{]}

\emph{{[}Jump condition at \(t=0\): reading \(F(s)\) as the impulse
response of the zero-state 5th-order ODE with denominator \((s-2)^5\)
(relative degree \(4\)). Phase vector
\(\mathbf y=(y,y',y'',y''',y^{(4)})\) changes by
\(\Delta\mathbf y(0)=(0,0,0,1,10)\).{]}}

Example 6.12 \[
\mathcal{L}^{-1} \left\{ \frac{1 + e^{-3s}}{s^4} \right\} \implies \frac{1}{6} \left[ t^3 + (t-3)^3 u(t-3) \right]
\] \emph{{[}WolframAlpha:}

\begin{lstlisting}
InverseLaplaceTransform[(1 + Exp[-3 s])/s^4, s, t]
\end{lstlisting}

\emph{Returns
\(f(t) = \frac{1}{6}\left[t^3 + (t-3)^3 \, u(t-3)\right]\)}{]}

\emph{{[}Jump conditions: the numerator
\(1+e^{-3s}=\mathcal L\{\delta(t)+\delta(t-3)\}\), so this is the
zero-state response of the 4th-order pure-integrator ODE \(y''''=x(t)\)
to two unit impulses. Phase vector \(\mathbf y=(y,y',y'',y''')\) changes
by \(\Delta\mathbf y(0)=(0,0,0,1)\) at \(t=0\), and by
\(\Delta\mathbf y(3)=(0,0,0,1)\) at \(t=3\).{]}}

Example 6.13 \[
\mathcal{L}^{-1} \left\{ \frac{s}{(s^2 + 4)^2} \right\} \implies \frac{1}{4} t \sin 2t
\] \emph{{[}WolframAlpha:}

\begin{lstlisting}
InverseLaplaceTransform[s/(s^2 + 4)^2, s, t]
\end{lstlisting}

\emph{Returns \(f(t) = \frac{1}{4} t \sin(2t)\)}{]}

\emph{{[}Jump condition at \(t=0\): reading \(F(s)\) as the impulse
response of the zero-state 4th-order ODE with denominator \((s^2+4)^2\)
(relative degree \(3\)). Phase vector \(\mathbf y=(y,y',y'',y''')\)
changes by \(\Delta\mathbf y(0)=(0,0,1,0)\).{]}}

Example 6.14 \[
\mathcal{L}^{-1} \left\{ \frac{8}{(s-1)(s^2+2s+5)} \right\} \implies f(t) = \mathcal{L}^{-1}\{F(s)\} = e^t - e^{-t} \cos 2t - e^{-t} \sin 2t
\] \emph{{[}WolframAlpha:}

\begin{lstlisting}
InverseLaplaceTransform[8/((s - 1)(s^2 + 2 s + 5)), s, t]
\end{lstlisting}

\emph{Returns \(f(t) = e^{t} - e^{-t}\cos(2t) - e^{-t}\sin(2t)\)}{]}

\emph{{[}Jump condition at \(t=0\): reading \(F(s)\) as the impulse
response of the zero-state 3rd-order ODE with denominator
\((s-1)(s^2+2s+5)\) (relative degree \(3\)). Phase vector
\(\mathbf y=(y,y',y'')\) changes by \(\Delta\mathbf y(0)=(0,0,8)\).{]}}

Example 6.15 \[
\mathcal{L}^{-1} \left\{ \frac{s+1}{(s^2+1)(s^2+9)} \right\} \implies f(t) = \mathcal{L}^{-1}\{F(s)\} = \frac{1}{8} \left( \cos t + \sin t - \cos 3t - \frac{1}{3} \sin 3t \right)
\] \emph{{[}WolframAlpha:}

\begin{lstlisting}
InverseLaplaceTransform[(s + 1)/((s^2 + 1)(s^2 + 9)), s, t]
\end{lstlisting}

\emph{Returns
\(f(t) = \frac{1}{8}\left(\cos t + \sin t - \cos 3t - \frac{1}{3}\sin 3t\right)\)}{]}

\emph{{[}Jump condition at \(t=0\): reading \(F(s)\) as the impulse
response of the zero-state 4th-order ODE with denominator
\((s^2+1)(s^2+9)\) (relative degree \(3\)). Phase vector
\(\mathbf y=(y,y',y'',y''')\) changes by
\(\Delta\mathbf y(0)=(0,0,1,1)\).{]}}

\begin{center}\rule{0.5\linewidth}{0.5pt}\end{center}

\hypertarget{additional-exercises}{%
\section{Additional Exercises}\label{additional-exercises}}

Supplementary problems and variations, given as page/problem references
only --- no solution is reproduced here. Sorted alphabetically by
author; intended to guide further practice.

\hypertarget{bottega}{%
\subsection*{Bottega}\label{bottega}}

p.238: Example 4.2 ``A tethered 1 pound ball hangs in the vertical plane
when it is tapped with a racket. Following the tap the ball is observed
to exhibit oscillatory motion of amplitude 0.2 radians with a period of
2 seconds. Determine the impulse imparted by the racket.'' (Problem
statement only --- no closed-form solution given in the source, hence
listed here rather than under Solved Exercises.)

p.236-238 couple 2nd order system with impulse load; p.269: Ex.4.4-4.6;
p.429: MDOF system under impulse load; p.470: double pendulum under
impulse load; p.488: Ex 8.17 elastically supported frame under struck;
p.501 Ex 8.7; p.507: Ex 8.26; p.715: Ex 11.3 PDE ``Determine the
response of the rod it is struck on its right end by an impulse of
magnitude''; p.718: Ex 11.17 ``The beam is impacted at its left end''

\hypertarget{boyce-diprima}{%
\subsection*{Boyce \& DiPrima}\label{boyce-diprima}}

p.273-274 has many exercises

\hypertarget{campbell-haberman}{%
\subsection*{Campbell \& Haberman}\label{campbell-haberman}}

p.264 Exercises 1--8

\hypertarget{de-oliveira}{%
\subsection*{De Oliveira}\label{de-oliveira}}

p.82 Problem 3.41 ``Compute the inverse Laplace transform of the
following complex-valued functions''; p.83 Problem 3.44 ``Compute the
inverse Laplace transform''

\hypertarget{dorf-bishop}{%
\subsection*{Dorf \& Bishop}\label{dorf-bishop}}

p.174 P2.36 ``Determine the impulse response of the system''; p.178
``Consider the unity feedback system described in the block diagram
\ldots{} Compute analytically the response of the system to an impulse
disturbance''; p.392 CP5.1 ``Obtain the impulse response analytically''

\hypertarget{edwards-penney}{%
\subsection*{Edwards \& Penney}\label{edwards-penney}}

9.326 4.6 Problems 1-8, 15-16 (equality of solution by changing IC)

\hypertarget{esfandiari-lu}{%
\subsection*{Esfandiari \& Lu}\label{esfandiari-lu}}

p.59 Problems 19 through 24; p.62 Problem 10 ``Solve the IVP''; p.352
``8.3.2 Impulse Response of Second-Order Systems''; p.353 ``Example 8.5:
Impulse Response''; p.359 ``Example 8.8: Impulse Response''; p.363
``PROBLEM SET 8.2''/7-11, 20

\hypertarget{franklin-powell-emami-naeini}{%
\subsection*{Franklin, Powell \&
Emami-Naeini}\label{franklin-powell-emami-naeini}}

p.230 EXAMPLE 4.9; p.589 Problem 7.20

\hypertarget{inman}{%
\subsection*{Inman}\label{inman}}

p.221 Example 3.1.1; p.222 Example 3.1.3; p.224 Example 3.1.4; p.232
Example 3.2.3; p.245 Example 3.4.4, Example 3.4.5; p.287 Problems
3.1-3.6, 3.10-3.13; p.377 Example 4.8.1 MDOF system with impulse; p.382
Example 4.8.2 MDOF system with impulse; p.386 Example 4.8.3 MDOF system
with impulse; p.428 Problem 4.76; p.429 Problem 4.78; p.440 Example
5.1.2; p.557 Example 6.8.1 ``Calculate the forced response of the string
fixed at both ends \ldots{} subject to unit impulse''; p.571 Problem
6.67

\hypertarget{karris}{%
\subsection*{Karris}\label{karris}}

p.6-2 ``Example 6.1''; p.6-3 ``Example 6.2''

\hypertarget{kelly}{%
\subsection*{Kelly}\label{kelly}}

p.317 EXAMPLE 5.1; p.374 Problem 5.21-5.23

\hypertarget{lathi-green}{%
\subsection*{Lathi \& Green}\label{lathi-green}}

p.471 Problem 4.3-6

\hypertarget{meirovitch}{%
\subsection*{Meirovitch}\label{meirovitch}}

p.371 Problem 7.49; p.463 Problem 8.38, 8.42, 8.44

\hypertarget{nagle-saff-snider}{%
\subsection*{Nagle, Saff \& Snider}\label{nagle-saff-snider}}

p.404 ``7.8 EXERCISES'' 5-12, 23-28; p.410 ``7.9 Exercises'' 13-29, 35;
p.416 ``REVIEW PROBLEMS FOR CHAPTER 7'' Problem 29-30

\hypertarget{ogata}{%
\subsection*{Ogata}\label{ogata}}

p.196 MATLAB Program 5--8 ``Unit-Impulse Response of G(s) = 1/(s\^{}2 +
0.2s + 1)''; p.264 B--5--4 ``Consider the system shown in Figure 5--72.
The system is initially at rest. Suppose that the cart is set into
motion by an impulsive force whose strength is unity. Can it be stopped
by another such impulsive force?''; p.264 B-5-5, B-5-6; p.265 B-5-10/11;
p.267 B-5-16

\hypertarget{rao}{%
\subsection*{Rao}\label{rao}}

p.382 ``4.5.1 Response to an Impulse''; p.384 EXAMPLE 4.7 ``Response of
a Structure Under Impact''; p.385 EXAMPLE 4.8 ``Response of a Structure
Under Double Impact''; p.407 EXAMPLE 4.9 ``Unit Impulse Response of a
First-Order System''; p.409 EXAMPLE 4.21 ``Unit Impulse Response of a
Second-Order System''; p.437 EXAMPLE 4.33 ``Impulse Response of a
Structure''; p.511 EXAMPLE 5.12 ``Response Under Impulse Using Laplace
Transform Method''

\hypertarget{schiff}{%
\subsection*{Schiff}\label{schiff}}

p.87 Exercises 2.5 1-7

\hypertarget{shabana}{%
\subsection*{Shabana}\label{shabana}}

p.45 Problems 1.3, 1.9

\hypertarget{thorby}{%
\subsection*{Thorby}\label{thorby}}

p.51 Example 3.2

\hypertarget{xue-chen-atherton}{%
\subsection*{Xue, Chen \& Atherton}\label{xue-chen-atherton}}

p.76 Example 3.20. ``Consider again the system model studied in Example
3.17. The impulse response of the system can be obtained as shown in
Figure 3.11:''
\passthrough{\lstinline!>> G=tf([10 20],[10 23 26 23 10],'ioDelay',1); impulse(G, 30);!}

p.106 Problem 9 ``Find impulse response for the system:'' \[
\frac{18s^7 + 514s^6 + 5982s^5 + 36380s^4 + 122664s^3 + 222088s^2 + 185760s + 40320}{s^8 + 36s^7 + 546s^6 + 4536s^5 + 22449s^4 + 67284s^3 + 118124s^2 + 109584s + 40320}
\]

\begin{center}\rule{0.5\linewidth}{0.5pt}\end{center}

\hypertarget{verification-methodology}{%
\section{Verification Methodology}\label{verification-methodology}}

Unlike a plain literature survey, every one of the thirty-seven solved
exercises in this compilation has been independently checked against its
own stated differential equation (or transfer function) and initial
conditions, rather than simply quoted. Three complementary methods were
used, chosen per entry according to what its stated form allowed:

\begin{itemize}
\tightlist
\item
  \textbf{Symbolic Laplace-transform matching.} For entries given as a
  transfer function \(H(s)\) (or \(C(s)\)) paired with an impulse/step
  response \(h(t)\) (or \(c(t)\)), and for IVPs whose claimed solution
  has no time-shifted (Heaviside) piece, the claimed time-domain
  solution's Laplace transform was computed symbolically (via SymPy) and
  compared against \(H(s)\), or against \(Y(s)\) built from the stated
  ODE's coefficients and initial conditions.
\item
  \textbf{Segment-wise ODE and jump-matching.} For IVPs whose forcing
  includes one or more shifted delta impulses \(\delta(t-c)\) --- where
  the claimed solution is naturally piecewise/Heaviside-driven --- each
  solution was checked directly: the homogeneous differential equation
  on every open interval between impulses, continuity of
  \(y, y', \dots, y^{(n-2)}\) at each impulse location, and the required
  jump of \(y^{(n-1)}\) by (impulse amplitude)/(leading coefficient) at
  that point. This avoids a limitation of general-purpose symbolic
  Laplace-transform routines, which do not reliably transform
  expressions built from
  \passthrough{\lstinline!Heaviside(t-c)*f(t-c)!}.
\item
  \textbf{Numeric self-consistency.} For the one entry specified purely
  by numeric physical parameters (Shabana, Example 1.10), the stated
  parameters were used to recompute the natural frequency, damping
  ratio, damped frequency, and response amplitude and check them against
  every numeric value the source itself reports.
\end{itemize}

This process confirmed all thirty-seven solved exercises exactly as
transcribed. It also surfaced four transcription issues beyond the
duplicate example discussed in the Conclusion below, each corrected
inline with an explicit derivation rather than silently: an incomplete
equation (Nagle, p.403, Example 4), a garbled unit label (Shabana, p.41,
Example 1.10), a single-character transcription error (Xie, p.258,
Example 6.13), and one example whose printed answer did not actually
satisfy its own stated equation and initial conditions (Gangadharaiah \&
Sandeep, p.340, Example 4.25).

A machine-readable export of the full, verified example set --- problem
source, coefficients or transfer function, forcing, initial conditions,
closed-form solution, and verification method for each of the
thirty-seven entries --- is provided alongside this article as
\passthrough{\lstinline!solved\_examples.json!} and
\passthrough{\lstinline!solved\_examples.csv!}, so the collection can be
consumed directly by an automated test harness rather than
re-transcribed by hand.

Concretely, each entry in
\passthrough{\lstinline!solved\_examples.json!} carries a
\passthrough{\lstinline!type!} field that determines how it is best
exercised as a test case:

\begin{itemize}
\tightlist
\item
  \textbf{\passthrough{\lstinline!ivp!} entries} (12 of the 37) give
  fully numeric ODE coefficients, a list of forcing impulses (amplitude,
  derivative order, and shift location), and numeric initial conditions
  --- a direct, ready-to-parse input for an ODE solver under test. A
  test harness can feed
  \passthrough{\lstinline!(ode\_coeffs\_highest\_to\_lowest, forcing, initial\_conditions)!}
  straight into the solver, evaluate both the solver's output and the
  quoted \passthrough{\lstinline!solution\_latex!} (parsed via a CAS
  such as SymPy) at a grid of sample times away from the impulse
  locations, and assert numerical agreement to within a chosen
  tolerance.
\item
  \textbf{\passthrough{\lstinline!transfer\_function!} entries} (14 of
  the 37) pair a transfer function \passthrough{\lstinline!H\_s!} with
  its known impulse response \passthrough{\lstinline!h\_t!} --- suited
  to testing a Laplace-domain toolbox: run the transform under test on
  \passthrough{\lstinline!H\_s!} and diff the result against
  \passthrough{\lstinline!h\_t!}.
\item
  \textbf{\passthrough{\lstinline!ivp\_symbolic!} entries} (10 of the
  37) carry general, symbolic parameters (e.g.~\(\tau\), \(k\),
  \(\omega_0\), \(\zeta\)) rather than fixed numbers, which makes them
  well suited to property-based or randomized testing: substitute random
  concrete values for the symbolic parameters before each comparison,
  exercising the solver across a swept parameter range instead of one
  fixed case.
\item
  \textbf{The one \passthrough{\lstinline!ivp\_numeric!} entry}
  (Shabana, Example 1.10) is a fully numeric physical-parameter case,
  suited to a straightforward fixed-input regression test.
\end{itemize}

Across all four types, the \passthrough{\lstinline!verified!} and
\passthrough{\lstinline!verification\_method!} fields let a harness
filter to only independently-checked entries before trusting them as an
oracle (all 37 here are \passthrough{\lstinline!verified: true!}), and
the stable \passthrough{\lstinline!id!} field gives each entry a natural
key for parametrized test naming
(e.g.~\passthrough{\lstinline!pytest.mark.parametrize!} keyed by
\passthrough{\lstinline!id!}, or a JUnit/xUnit test-case name). The flat
\passthrough{\lstinline!solved\_examples.csv!} carries the same
bibliographic and verification columns (\passthrough{\lstinline!id!},
\passthrough{\lstinline!author!}, \passthrough{\lstinline!book\_title!},
\passthrough{\lstinline!page!}, \passthrough{\lstinline!example!},
\passthrough{\lstinline!type!}, \passthrough{\lstinline!verified!},
\passthrough{\lstinline!verification\_method!},
\passthrough{\lstinline!notes!}) for quick spreadsheet review or for a
lightweight runner that would rather avoid a JSON parser; the fuller
mathematical content --- coefficients, forcing, initial conditions, and
closed-form solution --- is only in the JSON, since it does not flatten
cleanly into CSV columns.

\begin{center}\rule{0.5\linewidth}{0.5pt}\end{center}

\hypertarget{conclusion}{%
\section{Conclusion}\label{conclusion}}

This article gathered thirty-seven fully worked examples --- each an LTI
ODE (or, equivalently, a transfer function) forced by a Dirac delta
impulse or its derivative, paired with a published closed-form solution
--- from twelve textbooks spanning ordinary differential equations,
vibrations, signals and systems, and control engineering. These are
supplemented by dozens of further exercise references, without
solutions, indexed across twenty texts (several of which overlap with
the solved set) to point the reader toward additional practice material.
All content in both categories was extracted, not derived: no new
analytical results are claimed, and every solved example was
independently verified against its stated equation rather than merely
transcribed (see Verification Methodology above).

Three intended uses motivated the compilation. As a pedagogical
resource, the alphabetized, dual-category structure lets a reader move
directly from a specific author or problem type to the relevant worked
solution, without first locating and cross-referencing dozens of
separate books. As a software-engineering resource, the thirty-seven
solved exercises --- together with their machine-readable export and
per-example WolframAlpha check --- constitute a ready-made benchmark
suite: each pairs a well-posed initial value problem with an
independently verified analytical answer, suitable for regression or
unit testing of symbolic or numerical ODE solvers. As a small act of
literature quality control, the verification pass itself demonstrates
that even a well-regarded, widely used textbook can carry an internally
inconsistent worked answer (Gangadharaiah \& Sandeep, Example 4.25) or a
duplicated problem (Example 3.9, byte-identical to Example 3.12) that a
reader is unlikely to catch without redoing the algebra --- underscoring
the value of checking, not just collecting, textbook exercises before
relying on them for testing purposes.

Two limitations should be noted. First, the survey is not exhaustive: it
reflects the books available to the present compiler and is best
understood as a personal, growing reading list rather than a systematic
literature search. Second, while every solved example's closed-form
solution was checked against its own stated equation and initial
conditions, this verification cannot detect an error present identically
in both the stated problem and its stated answer (e.g., a genuine typo
in the source textbook's equation that happens to be consistent with its
own --- equally mistaken --- answer key); nor does it substitute for
tracing each problem back to first principles. One originally duplicated
example (Gangadharaiah \& Sandeep, Example 3.9) was removed rather than
repaired, since its true content could not be recovered independently of
Example 3.12. Extending the collection to further textbooks, and
extending the automated verification to the Additional Exercises once
solutions are added for them, are natural directions for future work.

\begin{center}\rule{0.5\linewidth}{0.5pt}\end{center}

\hypertarget{references}{%
\section{References}\label{references}}

Bottega, W. J. (2006). \emph{Engineering vibrations}. CRC Press. (ISBN:
9780849334207, 0849334209)

Boyce, W. E., \& DiPrima, R. C. (2017). Elementary differential
equations and boundary value problems (11th ed.). John Wiley \& Sons,
Inc.~(ISBN: 978-1-119-38164-8)

Campbell, S. L., \& Haberman, R. (2008). Introduction to differential
equations with dynamical systems. Princeton University Press. (ISBN:
978-0-691-12474-6)

De Oliveira, M. C. (2017). Fundamentals of linear control: A concise
approach. Cambridge University Press.
https://doi.org/10.1017/9781316941409 ISBN 978-1-107-18752-8 Hardback

Dorf, R. C., \& Bishop, R. H. (2008). Modern control systems: Solution
manual (11th ed.). Pearson Education, Inc.~(ISBN: 0-13-227029-3)

Edwards, C. H., \& Penney, D. E. (2008). Elementary differential
equations with boundary value problems (6th ed.). Pearson Education.
ISBN 0-13-600613-2

Esfandiari, R. S., \& Lu, B. (2014). Modeling and analysis of dynamic
systems (2nd ed.). CRC Press, Taylor \& Francis Group. (ISBN:
978-1-4665-7495-3) https://doi.org/10.1201/b16907

Franklin, G. F., Powell, J. D., \& Emami-Naeini, A. (2015). Feedback
control of dynamic systems (7th ed., Global ed.). Pearson Education
Limited. (ISBN: 978-1-292-06890-9)

Gangadharaiah, Y. H., \& Sandeep, N. (2021). Engineering applications of
the Laplace transform. Cambridge Scholars Publishing. ISBN (13):
978-1-5275-7373-4

Inman, D. J. (2014). Engineering vibration (4th ed.). Pearson Education,
Inc.~(ISBN: 978-0-13-287169-3)

Karris, S. T. (2003). Signals and systems with MATLAB® applications (2nd
ed.). Orchard Publications. (ISBN: 9780970951168, 0970951167)

Kelly, S. G. (2012). Mechanical vibrations: Theory and applications, SI.
Cengage Learning. (ISBN: 9781439062142)

Lathi, B. P., \& Green, R. A. (2018). Linear systems and signals (3rd
ed.). Oxford University Press. (ISBN: 978-0-19-020017-6)

Meirovitch, L. (2001). Fundamentals of vibrations (International ed.).
McGraw-Hill. (ISBN: 0-07-118174-1)

Nagle, R. K., Saff, E. B., \& Snider, A. D. (2018). Fundamentals of
differential equations (9th ed.). Pearson Education, Inc.~(ISBN:
978-0-321-97706-9)

Nagy, G. (n.d.). Ordinary differential equations. Mathematics
Department, Michigan State University

Ogata, K. (2010). Modern control engineering (5th ed.). Pearson
Education, Inc.~(ISBN- 13: 978-0-13-615673-4)

Rao, S. S. (2011). Mechanical vibrations (5th ed.). Pearson Education.
ISBN 978-0-13-212819-3

Schiff, Joel L. (1999). The Laplace transform: Theory and applications.
Springer-Verlag New York, Inc.~(ISBN: 0-387-98698-7)
https://doi.org/10.1007/978-0-387-22757-3

Shabana, A. A. (1997). Vibration of discrete and continuous systems (2nd
ed.). Springer-Verlag. https://doi.org/10.1007/978-1-4612-4036-5 Print
ISBN-13: 978-1-4612-8474-1

Thorby, D. (2008). Structural dynamics and vibration in practice: An
engineering handbook. Butterworth-Heinemann, an imprint of Elsevier.
(ISBN: 978-0-7506-8002-8)

Xie, W.-C. (2010). Differential equations for engineers. Cambridge
University Press. ISBN-13 978-0-521-19424-2

Xue, D., Chen, Y., \& Atherton, D. P. (2007). Linear feedback control:
Analysis and design with MATLAB. Society for Industrial and Applied
Mathematics. ISBN 978-0-898716-38-2

\end{document}
